\input{../tex-inputs/latex-leseansicht-vorspann}

\section[Arthur Schnitzler an Gustav Schwarzkopf, 12. 4. 1906]{L04040 Arthur Schnitzler an Gustav Schwarzkopf, 12. 4. 1906}
\nopagebreak\mylabel{L04040v}
\rehead{ }\normalsize\beginnumbering\briefempfaengerindex{Schwarzkopf, Gustav@\textsc{Schwarzkopf, Gustav}!zzzSchnitzler, Arthur@\emph{von Arthur Schnitzler}!1906-04-121@{12. 4. 1906}|(be}
\toendnotes[C]{\smallbreak\pagebreak[2]}
\correspDesc{Versand  durch Arthur Schnitzler am 12. 4. 1906 in Wien
\newline{}Übermittlung  am 13. 4. 1906 in Wien
\newline{}Erhalt  durch Gustav Schwarzkopf am 13. 4. 1906 in Wien}\toendnotes[C]{\smallbreak}
\Standort{CUL, Schnitzler, B 96.}
\physDesc{Postkarte, 353 Zeichen
\newline{}Handschrift: Bleistift, deutsche Kurrent
\newline{}Versand: 1) Rohrpost  2) Stempel: »\nobreak{}\oindex{XVIII., Währing@\textbf{XVIII., Währing}, \emph{Verwaltungsgebiet}|pwk}Wien 18/1 110, 13. IV. 06, 8–N\nobreak{}«.  3) Stempel: »\nobreak{}\oindex{I., Innere Stadt@\textbf{I., Innere Stadt}, \emph{Verwaltungsgebiet}|pwk}Wien 1/1, 23 IV 06, 8 20N\nobreak{}«. }\toendnotes[C]{\smallbreak}\pstart{}{\pb}Herr Gustav
                  Schwarzkopf\pend{}\pstart{}Wien I\oindex{I., Innere Stadt@\textbf{I., Innere Stadt}, \emph{Verwaltungsgebiet}|pw}\pend{}\pstart{}Tiefer Graben 23\oindex{Wien@\textbf{Wien}!I., Innere Stadt@\textbf{I., Innere Stadt}!Tiefer Graben 23@\textbf{Tiefer Graben 23}, \emph{Wohngebäude}|pw}.\pend{}{\bigskip}\vspace{1em}
\pstart
           {\pb}\textcolor{gray}{\textbf{Dr. Arthur Schnitzler}}\hfill 12/4. 906\pend
           
\pstart
           \textcolor{gray}{\textbf{Wien,
                        XVIII. Spoettelgasse 7\oindex{Wien@\textbf{Wien}!XVIII., Währing@\textbf{XVIII., Währing}!Edmund-Weiß-Gasse@\textbf{Edmund-Weiß-Gasse}, \emph{Straße}|pw}.}}\pend
           \vspace{0.5em}
\pstart
           lieber Guſtav, wollen Sie mit Bruder Max\pwindex{Schwarzkopf, Max 12.\,6.\,1857 Wien – 14.\,4.\,1928 ebd.@\textsc{Schwarzkopf, Max} (12.\,6.\,1857 Wien – 14.\,4.\,1928 ebd.), \emph{Rechtsanwalt}|pw}{ }\label{K_L04040-1v}\edtext{Oſterſo{\geminationn}tag}{\lemma{\textnormal{\emph{Ostersonntag}}}\Cendnote{\textnormal{Vgl. A. S.: \emph{Tagebuch}, 15. 4. 1906.
               }}}\label{K_L04040-1}{ }½ 10 Uhr Vm uns zum Spazierengehen \introOben{}(Weidlingbach\oindex{Weidlingbach@\textbf{Weidlingbach}|pw})\introOben{} abholen und da{\geminationn} gleich beide bei uns zu Mittag ſpeiſen? Wir würden
               uns ſehr freuen. Wenn ja, bedarfs keiner Antwort.\pend
           
\pstart
           Herzlichst{\\[\baselineskip]} Ihr{\\[\baselineskip]}\spacefill\mbox{A.}\pend
           \leftskip=0em{}
\pstart
           \noindent{}Iſt das Wetter \uline{ſcheußlich}, da{\geminationn} bitte zu Mittag, 1, ½ 2 bei uns zu ſein.\pend
           \selectlanguage{ngerman}\endnumbering\briefempfaengerindex{Schwarzkopf, Gustav@\textsc{Schwarzkopf, Gustav}!zzzSchnitzler, Arthur@\emph{von Arthur Schnitzler}!1906-04-121@{12. 4. 1906}|)be}\mylabel{L04040h}
\begin{anhang}
\end{anhang}\newcommand{\dateiname}{L04040}\newcommand{\titel}{Arthur Schnitzler an Gustav Schwarzkopf, 12. 4. 1906}\newcommand{\editorInnen}{Herausgegeben von Jahnke, SelmaMüller, Martin Anton}\input{../tex-inputs/latex-leseansicht-abspann}
